% from ega4-16-fr.tex
\oldpage[IV-4]{25}
\[
  \xymatrix{
    Y\ar[d]_-{\Delta_f}&X\ar[l]_-j\ar[d]^-{\Delta_{f\circ j}}\\
    Y\times_Z Y&X\times_Z X\ar[l]^-{j\times_Z j}
  }
\]
identifiant $X$ au produit des $(Y\times_Z Y)$-préschémas $Y$ et
$X\times_Z X$ (I, 5.3.7). Comme $j\times_Z j$ est une immersion, on déduit
donc de cette remarque et de (16.2.2) que si $\Delta_{Y/Z}^n$ et
$\Delta_{X/Z}^n$ désignent les voisinages infinitésimaux d'ordre $n$ de $Y$
et $X$ pour les immersions canoniques $\Delta_f$ et $\Delta_{f\circ j}$
respectivement, on a un diagramme
\[
  \xymatrix{
    \Delta_{Y/Z}^n\ar[d]&\Delta_{X/Z}^n\ar[l]\ar[d]\\
    Y\times_Z Y&X\times_Z X\ar[l]^-{j\times_Z j}
  }
\]
faisant de $\Delta_{X/Z}^n$ le produit des $(Y\times_Z Y)$-préschémas
$\Delta_{Y/Z}^n$ et $X\times_Z X$. On peut encore dire que
$\mathscr{P}_{X/Z}^n$ s'identifie au faisceau d'anneaux
$\mathscr{P}_{Y/Z}^n\otimes_{\mathcal{O}_{Y\times_Z Y}}
\mathcal{O}_{X\times_Z X}$. Mais on voit aussitôt (par exemple en se ramenant
au cas affine) que
\[
  \mathcal{O}_{X\times_Z X}=
  \mathcal{O}_{Y\times_Z Y}/
  (p_1^*(\mathscr{K})+p_2^*(\mathscr{K}))
  \mathcal{O}_{Y\times_Z Y}.
\]
Donc $\mathscr{P}_{X/Z}^n$ s'identifie au quotient de
$\mathscr{P}_{Y/Z}^n$ par l'Idéal engendré par l'image dans
$\mathscr{P}_{Y/Z}^n$ de $p_1^*(\mathscr{K})+p_2^*(\mathscr{K})$. Mais par
définition cet Idéal est aussi engendré par
$\varpi_1(\mathscr{K})+\varpi_2(\mathscr{K})$. C.Q.F.D.

\begin{corollary}[16.4.21]
\label{IV.16.4.21}
Soient $f:Y\to Z$ un morphisme, $j:X\to Y$ une immersion. On a une suite
exacte de $\mathcal{O}_X$-Modules quasi-cohérents
\[
  \xymatrix{
    \mathscr{N}_{X/Y}\ar[r]&j^*(\Omega_{Y/Z}^1)\ar[r]&
    \Omega_{X/Z}^1\ar[r]&0.
  }
  \tag{16.4.21.1}
  \label{IV.16.4.21.1}
\]
\end{corollary}

Lorsque $X$, $Y$ et $Z$ sont affines, on retrouve ainsi la suite exacte
(0, 20.5.12.1).

\begin{corollary}[16.4.22]
\label{IV.16.4.22}
Si $f:X\to S$ est un morphisme localement de présentation finie,
$\mathscr{P}_{X/S}^n$ et $\Omega_{X/S}^1$ sont des
$\mathcal{O}_X$-Modules quasi-cohérents de présentation finie.
\end{corollary}

On est aussitôt ramené au cas où $S=\operatorname{Spec}(A)$ est affine,
$X=\operatorname{Spec}(B)$, où
$B=A[T_1,\ldots,T_r]/\mathfrak{K}$, $\mathfrak{K}$ étant un idéal de type
fini de $C=A[T_1,\ldots,T_r]$. On applique alors (16.4.20) avec $Z=S$,
$Y=\operatorname{Spec}(C)$ et $\mathscr{K}=\widetilde{\mathfrak{K}}$. Alors
$j^*(\mathscr{P}_{Y/Z}^n)$ est un $\mathcal{O}_X$-Module libre de rang fini
(16.4.10) et l'hypothèse sur $\mathfrak{K}$ entraîne que
$j^*(\mathcal{O}_Y\cdot d_{Y/Z}^n(\mathscr{K}))$ engendre un
$\mathcal{O}_X$-Module quasi-cohérent de type fini; d'où la conclusion.

\begin{proposition}[16.4.23]
\label{IV.16.4.23}
Soient $X$, $Y$ deux $S$-préschémas, $Z=X\times_S Y$ leur produit,
$p:X\times_S Y\to X$ et $q:X\times_S Y\to Y$ les projections canoniques.
Alors l'homomorphisme canonique
\[
  p_{Z/X/S}\oplus q_{Z/Y/S}:
  p^*(\Omega_{X/S}^1)\oplus q^*(\Omega_{Y/S}^1)
  \longrightarrow\Omega_{(X\times_S Y)/S}^1
  \tag{16.4.23.1}
  \label{IV.16.4.23.1}
\]
est bijectif.
\end{proposition}

